\section{Introduction}\label{sec:intro}
\begin{figure}[t]
\centering\footnotesize
\setlength{\fboxsep}{3pt}
\fbox{\begin{minipage}{0.96\linewidth}
\textbf{Question} (BIRD, \texttt{european\_football\_2}): \emph{List the top 5 leagues in ascending order of the number of goals made in all seasons combined.}\\[2pt]
{\scriptsize\ttfamily
$Q_g$: SELECT l.name, SUM(m.home\_team\_goal) + SUM(m.away\_team\_goal) \dots\ GROUP BY l.name ORDER BY \dots\ LIMIT 5\\
$Q_p$: SELECT l.name, SUM(m.home\_team\_goal + m.away\_team\_goal) \dots\ GROUP BY l.name ORDER BY \dots\ LIMIT 5}\\[3pt]
\begin{minipage}[t]{0.49\linewidth}\raggedright
\textbf{(a) Certified witness at $\mathcal A_{1f}$} (excerpt)\\[1pt]
\begin{tabular}{@{}l@{\,}c@{\,}c@{}}
\toprule
League & home\_goal & away\_goal\\
\midrule
N.\ Eredivisie & 3 & 10\\
N.\ Eredivisie & \textcolor{red}{\texttt{NULL}} & 10\\
\bottomrule
\end{tabular}\\[2pt]
$Q_g\mapsto 23$, \ $Q_p\mapsto 13$: \emph{non-equivalent}. Found by CaliberGraph (0.2\,s), random fuzzing and ParSEval; satisfies the DDL, types and formats.
\end{minipage}\hfill
\begin{minipage}[t]{0.47\linewidth}\raggedright
\textbf{(b) Against the data profile}\\[1pt]
In BIRD, \texttt{home\_team\_goal} is \texttt{NULL} in \textbf{0 of 25,979} matches. On \texttt{NULL}-free data $\textsc{sum}(a){+}\textsc{sum}(b)=\textsc{sum}(a{+}b)$ for every group, so the queries compute the same totals; CaliberGraph-R and fuzzing-R find no $\mathcal A_R$ witness. Both annotation passes: \emph{implausible data}.
\end{minipage}
\end{minipage}}
\caption{A certified counterexample that says little about the question. Under DDL- and format-derived admissibility, three tools ``disprove'' an accepted prediction; the difference needs a played match without a goal count, which never occurs in BIRD. \V{attrImplMin}--\V{attrImplMax}\,\% of the sampled witnesses of every tool rest on such data.}\label{fig:example}
\end{figure}


Natural-language interfaces to databases are ranked by \emph{execution accuracy} (EX): a predicted SQL query counts as correct if it returns the gold query's result on the benchmark database~\cite{yu2018spider,li2023bird}. The test is cheap and reproducible, but it examines a single database. Queries that compute different things --- one counts patients, the other lab records; one keeps rows with a \texttt{NULL} score, the other drops them; one breaks ties by scan order, the other returns all tied rows --- often agree on that database, and every such agreement inflates the leaderboard.

Two lines of work search for databases on which an EX-correct prediction and the gold query disagree. Test generators execute both queries on additional databases, fuzzed around the original data~\cite{zhong2020testsuite} or generated for coverage of a query plan~\cite{chen2025parseval}. Verifiers search for a distinguishing database in a bounded model of SQL: SpotIt~\cite{klopfenstein2026spotit} and SpotIt+~\cite{tremante2026spotitplus} report that a large fraction of the BIRD predictions accepted by EX are not equivalent to the gold query and conclude that EX over-states accuracy.

\textbf{This paper asks what a distinguishing database actually proves.} The question is a data-engineering one: a counterexample is a database, and whether it says anything about the question depends on whether such data can occur. We answer it with a controlled experimental study of five disprover families on ten Text-to-SQL systems and \V{nExc} accepted BIRD predictions, replicated on thirteen further systems and on Spider, organised around six research questions:
\begin{itemize}
\item \textbf{RQ1} How many accepted predictions can each tool certify as non-equivalent when every tool is judged by the same certifier, tiers and budget?
\item \textbf{RQ2} What do certified witnesses show about the question: a prediction error, a gold error, an ambiguity, or nothing but implausible data --- and who can tell?
\item \textbf{RQ3} Can the benchmark's own data make witnesses meaningful, and which search finds such witnesses?
\item \textbf{RQ4} How much do the answers change EX and the rankings of systems?
\item \textbf{RQ5} Which design choices matter, and what does an audit cost?
\item \textbf{RQ6} Do the results hold on another benchmark and on unseen systems?
\end{itemize}

\textbf{Findings.} Under one SQLite certifier, released verifier counterexamples lose about half of their pairs at the DDL tier and 82--92\,\% at the format tier, and our SQLite-native disprover, CaliberGraph, certifies the most EX-accepted predictions at these tiers (\V{nCGLonef} of \V{nExc}). Yet almost none of these witnesses satisfies a data profile of the benchmark instance, and whether they are evidence depends on what the judge knows about the data. An LLM annotator with access to the benchmark database finds that \V{attrImplMin}--\V{attrImplMax}\,\% of every tool's witnesses separate the queries only on implausible data, such as a \texttt{NULL} goal count (Fig.~\ref{fig:example}). Two blind human annotators who saw only the witness almost never reached that verdict; shown facts about the benchmark data, they did so for \V{hvTwoImpConsOne} of the \V{hvTwoConsOne} schema-tier witnesses on which they agreed ($\kappa=\V{hvTwoKappaBin}$ between them), but agreed with the AI item by item only moderately ($\kappa=\V{hvTwoKappaAIbinMax}$, below the pre-registered 0.60); a repair with real values removes the difference for \V{xVanPct}\,\% of the witnesses. We therefore turn the benchmark database into an admissibility tier by \emph{data profiling} --- observed domains and nullability, unique column combinations, inclusion and value dependencies --- and search its real sub-databases first. At this tier the implausible share falls from \V{pCGImpl}\,\% to \V{pCGRImpl}\,\%, and CaliberGraph-R certifies \V{nCGRAtR} pairs against \V{nFZRAtR} for random fuzzing with the same data access, a lead in counts that replicates on thirteen unseen systems but not on Spider's small instances, where fuzzing's witnesses are more often meaningful. The effect on EX lies between \V{ptsErrRMin}--\V{ptsErrOnefMax} points per system (witnesses weighted by attributed error rates, a lower bound) and 15--20 points (every schema-tier witness counted), and most meaningful witnesses expose latent defects of gold queries or questions.

\textbf{Contributions.}
\begin{enumerate}
\item \emph{An evaluation protocol for disprovers} (Section~\ref{sec:protocol}): an independent certifier on the benchmark's engine, graded admissibility tiers, budget-matched baselines, and tool-blind attribution of witness meaning, checked against a second model, two rounds of human annotation and an annotator-free repair test. Comparisons were pre-registered in three stages with code freezes; failed hypotheses are reported.
\item \emph{A data-profile admissibility tier} $\mathcal A_R$ (Section~\ref{sec:profile}) that turns instance-discovered constraints into checkable realism conditions for counterexample databases, applicable to the output of any tool, with its cost and coverage on BIRD.
\item \emph{CaliberGraph and CaliberGraph-R} (Section~\ref{sec:method}), disprovers that build candidates along the join graph and, for $\mathcal A_R$, retrieve real sub-databases on which the two queries' predicates disagree. We report which components matter and which do not.
\item \emph{An experimental analysis} (Section~\ref{sec:results}) on BIRD, thirteen unseen systems and Spider, with findings and recommendations for benchmark maintainers and tool builders (Table~\ref{tab:findings}), and a released list of \V{nIssues} BIRD questions whose gold query or wording our witnesses call into question, which overlaps only weakly with corrected annotations of the same benchmark~\cite{jin2026annotation}.
\end{enumerate}
All code, data, witnesses, annotations and pre-registrations are released (see Artifact Availability).
